%%% CHANGELOG
%%% Patch 0606a (Session 146, 2026-05-27)
%%% Umbrella registration of THEO-DSL-8: Vertex-Aligned O(delta^2) Closure Theorem.
%%%
%%% Combines into a single citable theorem package:
%%%   - Patch 0606 (vertex-aligned first-order structural, hypothesis H1_V)
%%%   - Patch 0591 (vertex-aligned second-order alpha_2 = -9/phi^2)
%%%
%%% Direct mirror of Patch 0605a (THEO-DSL-7 edge-aligned umbrella).
%%% Completes the 2x2 vertex/edge x first/second-order coefficient phase
%%% diagram at O(delta^2). Together with THEO-DSL-6 (dimensional growth),
%%% THEO-DSL-7 (edge-aligned closure), and THEO-DSL-8 (this work), the
%%% O(delta^2) orientation phase diagram is structurally complete.
%%%
%%% Rigor level: sketch-document Layer 3 (matches constituents).
%%% Single-artifact discipline: 22-for-22 since Patch 0585.

\documentclass[11pt]{article}
\usepackage{amsmath,amssymb,amsthm}
\usepackage{geometry}
\geometry{margin=1in}

\newtheorem{theorem}{Theorem}
\newtheorem{corollary}{Corollary}
\newtheorem{remark}{Remark}

\title{THEO-DSL-8: Vertex-Aligned $\mathcal{O}(\delta^2)$ Closure Theorem \\[2pt]
\large Umbrella Registration (Patches 0591 + 0606)}
\author{Hyperphysics Institute --- Conscious Point Physics Programme}
\date{Patch 0606a \quad Session 146 \quad 2026-05-27}

\begin{document}
\maketitle

\begin{abstract}
This document registers THEO-DSL-8 as the umbrella theorem covering the
vertex-aligned current expansion through second order in $\delta$. The
combined statement, drawn from Patch 0606 (first order) and Patch 0591
(second order), asserts that at vertex-aligned orientation $\theta_V$,
the $D_5$-invariant current $\mathbf{j}(\theta_V, \delta)$ is confined
order-by-order to the 1D substrate-primitive subspace
$\mathrm{span}\{\hat{\mathbf{n}}_{\text{subst}}\}$ through
$\mathcal{O}(\delta^2)$, with explicit, real, non-vanishing coefficients
$\alpha_1 = +6/\varphi^2$ and $\alpha_2 = -9/\varphi^2$ exhibiting a sign
inversion between orders. Together with THEO-DSL-7 (edge-aligned
$\mathcal{O}(\delta^2)$ closure, Patch 0605a) and THEO-DSL-6
(dimensional-growth structural), THEO-DSL-8 completes the
$\mathcal{O}(\delta^2)$ orientation phase diagram. Five umbrella-level
findings (F-DSL-8-1 through F-DSL-8-5) record cross-patch and
cross-umbrella observations.
\end{abstract}

\section{Theorem Statement (THEO-DSL-8)}

\begin{theorem}[Vertex-Aligned $\mathcal{O}(\delta^2)$ Closure --- THEO-DSL-8]
\label{thm:THEO-DSL-8}
Let $\mathbf{j}(\theta_V, \delta)$ denote the $D_5$-invariant current at
vertex-aligned orientation, expanded as
\[
\mathbf{j}(\theta_V, \delta) \;=\; \mathbf{j}_1\,\delta \;+\; \mathbf{j}_2\,\delta^2
\;+\; \mathcal{O}(\delta^3).
\]
Then through $\mathcal{O}(\delta^2)$:
\begin{enumerate}
\item[\textnormal{(i)}] \textbf{1D confinement.} Both $\mathbf{j}_1$ and
$\mathbf{j}_2$ lie within the 1D substrate-primitive subspace
$\mathrm{span}\{\hat{\mathbf{n}}_{\text{subst}}\}$, with the scalar
decompositions
\begin{align}
\mathbf{j}_1 &\;=\; \alpha_1\,\hat{\mathbf{n}}_{\text{subst}}, \\
\mathbf{j}_2 &\;=\; \alpha_2\,\hat{\mathbf{n}}_{\text{subst}},
\end{align}
establishing hypothesis (H1\_V) from Patch 0606 and its
vertex-aligned second-order analog from Patch 0591.

\item[\textnormal{(ii)}] \textbf{First-order coefficient} (Patch 0606
Theorem 1):
\[
\alpha_1 \;=\; +\frac{6}{\varphi^2} \;=\; 6(2-\varphi) \;=\; 12 - 6\varphi
\;\approx\; +2.292.
\]

\item[\textnormal{(iii)}] \textbf{Second-order coefficient} (Patch 0591
Theorem 5.1):
\[
\alpha_2 \;=\; -\frac{9}{\varphi^2} \;=\; -9(2-\varphi) \;=\; 9\varphi - 18
\;\approx\; -3.438.
\]

\item[\textnormal{(iv)}] \textbf{Closed-form coefficient ratio:}
\[
\frac{\alpha_2}{\alpha_1} \;=\; -\frac{3}{2},
\]
exactly and independent of $\varphi$ (the $\varphi^{-2}$ factors cancel).
The negative ratio is the vertex-aligned signature, in explicit contrast
to the edge-aligned positive ratios
$\alpha_{2,\rho}/\alpha_{1,\rho},\,
\alpha_{2,\text{edge}}/\alpha_{1,\text{edge}} > 0$ (Patch 0605
Finding A5-F4).
\end{enumerate}
\end{theorem}

\medskip
The rigor of THEO-DSL-8 inherits from its constituents: sketch-document
Layer 3 per DG-F15A-5 default-(A). No new proof obligations are
introduced by the umbrella; this registration provides a single citable
theorem statement for external reference.

\section{Constituent Patches and Results}

\subsection*{Patch 0606 --- Vertex-Aligned First-Order Structural}
\begin{itemize}
\item Hypothesis (H1\_V): $\mathbf{j}_1 \in
\mathrm{span}\{\hat{\mathbf{n}}_{\text{subst}}\}$ (1D collapse).
\item Lemma 1: 1D collapse at vertex-aligned via $D_5$ representation
branching (the $\mathcal{O}(\delta)$ realization of THEO-DSL-6
dimensional-growth).
\item Coefficient $\alpha_1 = +6/\varphi^2$ established.
\item Magnitude ratio $|\alpha_1^V|/|\alpha_1^{E,\rho}| = 2\varphi$
exactly (Finding A1V-F2).
\item Findings A1V-F1 through A1V-F4.
\item File:
\texttt{hardened\_theorems/vertex\_aligned\_invariant\_subspace\_structural.tex}
\end{itemize}

\subsection*{Patch 0591 --- Vertex-Aligned Second-Order $\alpha_2$}
\begin{itemize}
\item Vertex-aligned second-order coefficient
$\alpha_2 = -9/\varphi^2$ established.
\item Path-class enumeration §3: four $D_5$ orbits with weights
$\{1,\, 1/2,\, -1/(2\varphi),\, -\varphi/2\}$ (orientation-independent
per F-DSL-7-2).
\item Sub-class structure §4: $B_2 \equiv 0$, $B_4 = -\varphi/2$ uniform.
\item File:
\texttt{hardened\_theorems/o\_delta\_squared\_path\_class\_weights.tex}
\end{itemize}

\section{Umbrella-Level Findings}

These findings record cross-patch and cross-umbrella observations not
isolated at the level of either constituent alone.

\paragraph{F-DSL-8-1: Order-by-Order 1D Confinement.}
The same 1D substrate-primitive subspace
$\mathrm{span}\{\hat{\mathbf{n}}_{\text{subst}}\}$ confines the current
at \emph{both} $\mathcal{O}(\delta)$ and $\mathcal{O}(\delta^2)$. A priori,
second-order corrections could escape the first-order invariant subspace
via the symmetric tensor square; that they do not at vertex-aligned
reflects the rigidity of the vertex-stabilizer subgroup of the $D_5$
point group. This is the vertex-aligned analog of F-DSL-7-1
(where the surviving subspace at edge-aligned is 2D).

\paragraph{F-DSL-8-2: Vertex-Aligned Sub-Class Simplifications Persist Through $\mathcal{O}(\delta^2)$.}
The simplifications $B_2(O_i) \equiv 0$ and $B_4(O_i) = -\varphi/2$
(uniform) hold at vertex-aligned orientation at \emph{both} first
(Patch 0606) and second (Patch 0591) order. The persistence across
orders is a structural feature of the vertex-stabilizer. The
edge-aligned counterparts break uniformity at second order (Patch 0605
Remark~4.1 for $B_2$ non-vanishing; Patch 0605 Corollary~3.1 for $B_4$
splitting into four distinct values).

\paragraph{F-DSL-8-3: Sign Inversion as Vertex-Aligned Signature.}
$\alpha_1 = +6/\varphi^2 > 0$ at first order while
$\alpha_2 = -9/\varphi^2 < 0$ at second order: the coefficients flip
sign between orders. At edge-aligned, both first-order coefficients (per
Patch 0600) and both second-order coefficients (Patch 0605:
$\alpha_{2,\rho}, \alpha_{2,\text{edge}} > 0$) are non-negative; no sign
flip occurs (THEO-DSL-7 Theorem (iii)). The sign-flip / no-sign-flip
distinction is a candidate diagnostic for orientation phase
identification at higher orders.

\paragraph{F-DSL-8-4: Closed-Form Coefficient Ratio $\alpha_2/\alpha_1 = -3/2$.}
The vertex-aligned ratio $\alpha_2/\alpha_1$ evaluates exactly to $-3/2$,
independent of $\varphi$ because the $\varphi^{-2}$ factors cancel. This
contrasts the edge-aligned ratios
$\alpha_{2,\rho}/\alpha_{1,\rho}$ and
$\alpha_{2,\text{edge}}/\alpha_{1,\text{edge}}$, both positive but
$\varphi$-dependent (Patch 0605 Finding A5-F4).

\paragraph{F-DSL-8-5: $\mathcal{O}(\delta^2)$ Orientation Phase Diagram Complete.}
With THEO-DSL-7 (edge-aligned) and THEO-DSL-8 (vertex-aligned), the
$2 \times 2$ vertex/edge $\times$ first/second-order phase diagram of
the current expansion at $\mathcal{O}(\delta^2)$ is closed:

\begin{center}
\begin{tabular}{c|c|c}
\hline
 & Vertex-aligned (THEO-DSL-8) & Edge-aligned (THEO-DSL-7) \\
\hline
$\mathcal{O}(\delta)$
 & 1D, $\alpha_1 = +6/\varphi^2$
 & 2D, $(\alpha_{1,\rho}, \alpha_{1,\text{edge}})$ \\[2pt]
$\mathcal{O}(\delta^2)$
 & 1D, $\alpha_2 = -9/\varphi^2$
 & 2D, $(\alpha_{2,\rho}, \alpha_{2,\text{edge}})$ \\
 &
 & $= (+\tfrac{9}{2\varphi},\, +9\varphi-12)$ \\[2pt]
ratio $\alpha_2/\alpha_1$
 & $-3/2$ (closed form)
 & both positive, $\varphi$-dependent \\
\hline
\end{tabular}
\end{center}

Together with THEO-DSL-6 (dimensional-growth structural), the triad
\{THEO-DSL-6, THEO-DSL-7, THEO-DSL-8\} forms a coherent governing set
for the $\mathcal{O}(\delta^2)$ orientation phenomena. The natural next
moves are $\mathcal{O}(\delta^3)$ extension (open at both orientations)
or a meta-umbrella combining the triad.

\section{Registration Record}

\begin{center}
\begin{tabular}{ll}
\hline
\textbf{Registry identifier:} & THEO-DSL-8 \\
\textbf{Series:} & THEO-DSL (theorem-level package registrations) \\
\textbf{Constituent patches:} & 0591, 0606 \\
\textbf{Hypotheses covered:} & (H1\_V) and the vertex-aligned \\
 & second-order analog from Patch 0591 \\
\textbf{Rigor level:} & sketch-document Layer 3 (matches constituents) \\
\textbf{Governance designation:} & DG-F15A-5 default-(A) \\
\textbf{Companion umbrella:} & THEO-DSL-7 (edge-aligned, Patch 0605a) \\
\textbf{Registration session:} & 146 \\
\textbf{Registration date:} & 2026-05-27 \\
\textbf{Single-artifact discipline:} & 22-for-22 since Patch 0585 \\
\hline
\end{tabular}
\end{center}

\section{Citation Form}

When citing this package from external papers, use:

\medskip
\noindent\texttt{[THEO-DSL-8]} \emph{Vertex-Aligned $\mathcal{O}(\delta^2)$
Closure Theorem.} CPP Programme internal registry, Session 146,
2026-05-27. Constituents: Patch 0606 (first-order structural) and
Patch 0591 (second-order closure).

\section{Forward Items}

\begin{itemize}
\item \textbf{OPEN-DSL-8-1}: Layer 4 algebraic hardening of the
constituent orbital decompositions (depends on OPEN-A1V-1 from Patch 0606
for the first-order side).
\item \textbf{OPEN-DSL-8-2}: $\mathcal{O}(\delta^3)$ extension at
vertex-aligned orientation (companion to OPEN-DSL-7-2 at edge-aligned).
\item \textbf{OPEN-DSL-8-3}: Meta-umbrella combining THEO-DSL-6
(dimensional growth), THEO-DSL-7 (edge-aligned closure), and THEO-DSL-8
(this work) into a single $\mathcal{O}(\delta^2)$ orientation phase
diagram theorem. Candidate identifier: THEO-DSL-9.
\item \textbf{OPEN-DSL-8-4}: Cross-orientation interpolation
$\alpha_n(\theta)$ between $\theta_V$ and $\theta_E$ at each order $n$,
exploiting the sign-flip / no-sign-flip distinction F-DSL-8-3
(companion to OPEN-DSL-7-3).
\end{itemize}

\end{document}
